\section{Deployment Codetask}

Codetask verfügt über drei verschiedene Umgebungen: eine Local-, eine Staging- und eine Prod-Umgebung. 
Jede dieser Umgebungen hat ihre eigene Compose-Datei. Für die Haupt-API, die Check-API, die AI-API und das Webframework sind eigene Dockerfiles angelegt, wobei jeweils für die lokale Umgebung ein eigenes Dockerfile verwendet wird. 
Jede dieser Umgebungen hat eine eigene .env-Datei, in der Passwörter, Tokens und weitere Variablen stehen, die im Code aufgerufen werden.

Im Ordner .deployment gibt es die Datei mongo-init.js, die verwendet wird, um die MongoDB zu initialisieren.
\begin{figure}[htbp]

\dirtree{%
.1 teamprojekt/.
.2 .deployment/.
.3 prod/.
.4 docker-compose.yml.
.3 staging/.
.4 docker-compose.yml.
.3 mongo-init.js.
.2 .docker/.
.2 ai-api/.
.2 api/.
.2 check-api/.
.2 web/.
.3 frontend/.
.2 docker-compose.yml.
}

\caption{Vereinfachtes Projektverzeichnis der Codetask-Anwendung}
\label{fig:projektstruktur}

\end{figure}

\FloatBarrier
\subsection{Lokale Umgebung}

Die Docker-Compose.yml-Datei für die lokale Umgebung liegt im Root-Verzeichnis des Projekts.
Sie dient dazu, dass Codetask lokal ausgeführt werden kann, um daran zu entwickeln und Veränderungen auszuprobieren, bevor diese gepusht werden.
Daher sind alle Services unter localhost erreichbar, bekommen aber vom Compose jeweils eigene Ports zugeordnet. Es werden keine weiteren Netzwerkeinstellungen eingerichtet.

Das Compose der lokalen Umgebung baut die Images mit den selbst erstellten Dockerfiles selbst und nimmt dabei das Root-Verzeichnis als Kontext.
Die Anwendungen in den Containern werden jeweils im DEV-Modus gestartet, wodurch das Webframework beispielsweise einen eigenen Webserver hat. Dadurch verbrauchen alle Container jedoch mehr Ressourcen als im Normalbetrieb.
Die Images von nginx, mongodb und mongo-express werden aus dem Internet gezogen. 

Alle Services haben die Option „restart: unless-stopped” hinzugefügt bekommen, wodurch Container, wenn sie ausfallen, direkt neu gestartet werden.

Die Services haben ein Env-File zugeordnet bekommen, aber manche haben direkt eigene Environment-Variablen im Compose. Das Besondere an der lokalen Umgebung ist, dass die Volumes der Konfiguration der Services im Container überschrieben werden, wodurch die Konfiguration außerhalb des Containers angepasst werden kann, während dieser läuft, und nach einem Neustart verwendet werden kann.
\subsection{Staging und Production Umgebung}

Die Docker-Compose-Dateien sind im Unterordner /.deployment zu finden.
In der Produktionsumgebung läuft die aktuelle Live-Version von Codetask, während die Staging-Umgebung einen Zwischenstand zwischen der lokalen und der Produktionsumgebung darstellt.
In dieser können Implementierungen und Entwicklungen in einer live-ähnlichen Umgebung getestet werden.

Die Compose-Dateien ähneln denen der lokalen Umgebung, haben dennoch deutliche Unterschiede.
So gibt es beispielsweise einen weiteren Service namens Traefik, der als Reverse Proxy benutzt wird. 
Die Umgebungen haben auch einen Eintrag des DNS-Servers der HTWG, wodurch sie auch von außen erreichbar sind. Die Staging-Umgebung ist hingegen nur über VPN erreichbar.
Der Host der Prod-Umgebung ist codetask.in.htwg-konstanz.de, der Host der Staging-Umgebung ist staging-codetask.in.htwg-konstanz.de. Wird der Host für die API oder AI-API verwendet, wird jeweils das Passende vorne dran gehängt, z. B. „api.codetask.in.htwg-konstanz.de”.
Daher haben beide Umgebungen jeweils drei URLs, wodurch drei Services von außen erreichbar sind. Die restlichen Services sind nur intern im Compose-Netzwerk erreichbar.
Bei der Produktionsumgebung haben alle drei URLs ein TLS-Zertifikat erhalten, wodurch jede mit HTTPS genutzt werden kann. Bei der Staging-Umgebung hat nur die URL für das Webframework ein TLS-Zertifikat erhalten.

Anders als bei der lokalen Umgebung werden die Images bei beiden Compose nicht mehr mit Dockerfiles selbst gebaut, sondern über die Container Registry vom GitLab-Repository von Codetask geholt.
Die Staging- und Produktionsumgebung haben auch ein anderes Dockerfile als die lokale Umgebung. Ein Unterschied ist beispielsweise, dass die Container nicht im DEV-Modus laufen oder dass im Frontend-Dockerfile ein Container für das Webframework und den Webserver enthalten ist.